从某种意义上讲，热力学的发展与人们制造和改进热机的努力是分不开的。最初，人们设想制造一种
不需要能量输入却能够持续工作的热机，但热力学第一定律告诉我们，这样的努力没有任何希望。退而求
其次，人们又希望得到效率为百分之百的热机，即把吸收的热量全部转化为机械功。从能量守恒的角度看，
这种设想并不直接违背热力学第一定律：能量输入与输出可以相等，但这还没有回答过程是否能够发生。
第二定律额外指出，循环装置不能以“只从单一热库取热并把它全部变成功”为唯一效果。
常见气体热机确有膨胀做功、压缩复原和向外放热的过程，但这只能帮助理解具体模型，不能证明一般定律。
例如液体表面缩小时也能对外做功，后面将研究这种不同的工质。

人们在追求尽可能高效的热机时，逐渐发现这一事实背后隐藏着一条普遍的热力学原理，这就是今天的
热力学第二定律。热机、可逆性以及熵，将是本讲讨论的主题。

\subsection*{热力学循环的效率}

热机在工作时循环运行：工质从某个状态出发，经历一系列过程后再次回到初始状态，完成一个热力学循环。
对在两个热库之间工作的普通热机，工质从高温热库吸热，还向低温热库放热。把整个循环
从外界吸收的总热量记为 $Q^{\rm in}$，把向外排出的总热量记为 $Q^{\rm out}$。由于循环结束后工质回到
初态，内能改变量为零。这里 $Q^{\rm in},Q^{\rm out}$ 分别是累计吸热、累计放热的正量，
也记为 $Q_h,Q_c$；净输入热量为二者之差，不沿用上一章单段过程中带符号的 $Q^{\rm in}=Q$。
热力学第一定律给出循环的净功
\begin{equation}
  \boxed{W^{\rm out}=Q^{\rm in}-Q^{\rm out}}.
\end{equation}
若工质仅以体积功与外界交换机械能，且过程可逆，则净功为 $\oint P\dd V$，等于 $P$--$V$ 图上
围成的有向面积，顺时针为正。对于电功、表面功等形式，须使用相应的功表达式，不能一概用 $P$--$V$ 面积。
只要排出的热量不为零，热机就没有把吸收的全部热量变成功。定义热机
效率为对外净功与吸热量的比值：
\begin{equation}
  \boxed{\eta=\frac{W^{\rm out}}{Q^{\rm in}}
  =1-\frac{Q^{\rm out}}{Q^{\rm in}}}.
\end{equation}

最常见的内燃机可用 Otto 热机近似描述。图21.1把活塞的工作过程与相应的 $P$--$V$ 循环放在一起。
从活塞运动看，一个循环可分为进气、压缩、点火、做功和排气几个环节。首先关闭排气口、打开进气口，
空气和汽油的混合气体被吸入气缸；自然吸气发动机的进气压强接近一个大气压，涡轮增压发动机则更高。
活塞到达最下端后进气口关闭，活塞因惯性向上运动并压缩混合气体。压缩过程通常很快，虽然气体温度
升高，但与气缸的热量交换可以忽略，因而可近似为绝热压缩。气体被完全压缩后，火花塞点燃高压油气
混合物，燃烧很快，活塞来不及移动，气体体积近似不变而温度、压强急剧升高。随后高温高压气体推动
活塞向下运动，进入做功冲程；这一快速膨胀过程也可近似为绝热膨胀。活塞到达最下端后排气口打开，
废气被排出，装置完成一个循环。

\ottoCycleFigure
\diagramcaption{图 21.1：上方是四冲程发动机的五个动作阶段，下方是空气标准奥托循环。点火单列显示，不是额外增加一个活塞冲程。}

最小、最大气缸体积分别为 $V_2,V_1$。先看真实进排气：$O\to A$ 提示吸气，压强近似不变，
但缸内气体量不断增加，并非定量气体等压膨胀；$A\to O$ 提示排气，气体不断离开气缸。
在真实做功冲程后打开排气门时，废气也会先因内外压差排出。这些开放过程伴随物质携能，
不能套用一份封闭气体的 $nc_v\Delta T$ 来计算排出能量。

为了求一个便于比较的理论效率，现在换用\textbf{空气标准模型}：用同一份定量、定组成的理想气体
代替不断更新的混合气；用可逆等容吸热代替燃烧，用可逆等容放热代替排气散热的效果。
蓝色闭式循环只有四段：$A\to B$ 绝热压缩，$B\to C$ 等容吸热，$C\to D$ 绝热膨胀，
$D\to A$ 等容放热。灰色 $O\leftrightarrow A$ 不是这份封闭工质实际走过的两段。

在这一理想化循环中，只有 $B\to C$ 的等容过程吸热。设气缸内气体可视为总物质的量为 $n$、定容
摩尔热容 $c_v$ 在整个温区内近似恒定的理想气体，并忽略摩擦、漏热和燃烧成分变化，则
\[
  Q_h=n c_v(T_C-T_B).
\]
类似地，只有 $D\to A$ 的等容过程向外排热，排热量为
\[
  Q_c=n c_v(T_D-T_A).
\]
所以
\[
  \eta=1-\frac{Q_c}{Q_h}
  =1-\frac{T_D-T_A}{T_C-T_B}.
\]
再利用两段绝热过程
\[
  T_A V_1^{\gamma-1}=T_B V_2^{\gamma-1},
  \qquad
  T_D V_1^{\gamma-1}=T_C V_2^{\gamma-1},
\]
令压缩比 $r=V_1/V_2>1$，两式分别给出 $T_B=r^{\gamma-1}T_A$、
$T_C=r^{\gamma-1}T_D$。相减而不是直接约去温度，可得
\[
 T_C-T_B=r^{\gamma-1}(T_D-T_A),\qquad
 \frac{Q_c}{Q_h}=\frac{T_D-T_A}{T_C-T_B}=r^{1-\gamma}.
\]
于是得到奥托循环效率与发动机几何参数的关系
\begin{equation}
  \boxed{\eta=1-\left(\frac{V_2}{V_1}\right)^{\gamma-1}}.
\end{equation}
这说明在理想模型中，压缩比 $V_1/V_2$ 越大，循环效率越高；真实发动机还会受到爆震、传热、摩擦和
材料强度等条件限制。

\subsection*{Carnot 循环}

在所有热力学循环中，对理论发展最重要的是 Carnot 设计的 Carnot 热机。它来自人们寻找最高可能效率
热机的努力；当理论上确定这种热机的效率以后，人们才认识到其背后隐藏着一条基本的热力学定律。
以下示范理想气体循环的计算时，仍取定量气体、恒定 $c_v,c_p$；卡诺定理本身并不限于这一模型。

\carnotCycleFigure
\diagramcaption{图 21.2：Carnot 循环的 $P$--$V$ 图。$1\to2$、$3\to4$ 为等温过程，$2\to3$、$4\to1$ 为绝热过程。}

如图21.2所示，Carnot 循环由两段绝热过程和两段等温过程组成，两条等温线的温度分别为 $T_1>T_2$。
正是两个热库之间存在温差，才有可能把一部分热能转化为机械功。循环从状态1出发：工质与温度为 $T_1$
的高温热库接触，持续吸热并作等温膨胀，工质对外做功；到达状态2后与热库脱离，继续绝热膨胀，直到
温度降到 $T_2$，来到状态3；然后工质与低温热库接触，受到外界做功而作等温压缩，同时向低温热库
排出热量 $Q_2$；到达状态4后再与热库脱离，作绝热压缩，最后恰好回到初态1。

Carnot 可以证明，无论采用何种工作物质，可逆 Carnot 循环都具有相同效率。先以理想气体为例。设热机
内有 $n$ 摩尔理想气体，高温等温膨胀吸入的热量与低温等温压缩排出的热量分别为
\[
  Q_1=nRT_1\ln\frac{V_2}{V_1},
  \qquad
  Q_2=nRT_2\ln\frac{V_3}{V_4}.
\]
两段绝热过程满足
\[
  T_1V_2^{\gamma-1}=T_2V_3^{\gamma-1},
  \qquad
  T_1V_1^{\gamma-1}=T_2V_4^{\gamma-1}.
\]
两式相除得到 $V_2/V_1=V_3/V_4$。于是效率表达式中所有体积因子恰好抵消，最终得到
\begin{equation}
  \boxed{\eta_{\rm Carnot}=1-\frac{T_2}{T_1}}.
\end{equation}
令人惊讶的是，Carnot 热机的效率只与两个热库的热力学温度有关，与热机的几何结构无关；这与 Otto
热机效率依赖压缩比形成鲜明对比。

\subsection*{热力学第二定律的早期陈述形式}

Carnot 关于热机效率的研究使人们逐渐意识到，有一条基本规律阻止人们制造效率为 $100\%$ 的热机。
这种机器并不直接违反热力学第一定律，却从未能够实现。图21.3(a)表示真实热机：它从高温热库吸收
$Q_h$，对外做功 $W=Q_h-Q_c$，并向低温热库排出 $Q_c$。图21.3(b)所示的“完美”热机则只从单一
热库取热并把它全部变成功。Lord Kelvin 彻底断绝了这种幻想，他指出：

\begin{quote}\small
不存在这样的装置：当它完成一个循环时，能从一个物体取出热量并把它转换为功，而不发生其他变化
（至少要在其他物体中发生变化）。
\end{quote}

历史上这种机器被称为第二类永动机，因此热力学第二定律的 Kelvin 表述也可理解为：第二类永动机不可能
制成。

\heatDevicesFigure
\diagramcaption{图 21.3：热机和制冷机能量流。由左至右为普通热机、第二类永动机、普通制冷机和无功制冷机；两种“完美”设想不允许实现。}

生活经验还指出，热量会自发地从高温物体传向低温物体。冰箱、空调试图把热量从低温部分搬运到高温
环境，必须依赖外部能量供应。在力学中，物体可以凭惯性短时间沿受力的反方向运动，例如竖直上抛物体
仍继续上升，或弹簧振子在回复力作用下仍远离平衡位置；不能据此类比出一种可持续的无补偿逆向传热。
在这里讨论的通常宏观传热中，若不付出其他代价，热量只能自发地从高温传向低温。Clausius 把这一事实表述为：

\begin{quote}\small
不存在这样的装置：当它完成一个循环时，能使热量从较低温度的热库传到较高温度的热库，而没有补偿
变化。
\end{quote}

从高温向低温的传热过程可以自然发生，而它的逆过程不能独立发生，因此 Clausius 表述也意味着：有限
温差下的热传导是一种不可逆过程。Kelvin 与 Clausius 的经典英文表述分别为：
\begin{quote}\footnotesize
It is impossible, by means of inanimate material agency, to derive mechanical effect from any portion of
matter by cooling it below the temperature of the coldest of the surrounding objects.

Heat can never pass from a colder to a warmer body without some other change, connected therewith,
occurring at the same time.
\end{quote}

\foundersPortraitFigure
\diagramcaption{图 21.4：三位热力学奠基人：卡诺研究热机的理想效率，开尔文和克劳修斯提出第二定律的重要表述。}
图21.4按卡诺、开尔文、克劳修斯的顺序排列，分别对应前述热机研究及两种早期陈述。

\subsection*{熵与热力学第二定律}

在得到热力学第二定律的早期陈述以后，人们逐渐认识到：对于任意热力学系统，除压强、密度、温度、
内能等状态变量外，还存在另一个状态变量——熵（entropy），记为 $S$。熵从出现之初就带有神秘色彩，
尤其是由熵增出发对宇宙结局的推测；正因为如此，更应准确理解熵的定义和适用条件，避免在不正确的
条件下滥用相关定律。

熵首先是状态变量。与内能一样，它完全由热力学系统的状态决定，而与系统以什么过程到达该状态无关。
状态变量必然同组成系统的微观粒子的性质有关：压强是粒子碰撞容器表面的宏观效果；在经典图像中，
温度由粒子的平均能量决定；内能可理解为全部微观粒子的能量总和。熵也由微观状态决定，Boltzmann 后来
揭示了这一联系。

当热力学系统经历某个过程时，熵可能改变。对于可逆过程中的任意状态，都可以定义系统温度。若系统在
温度为 $T$ 的平衡态吸收一小份热量 $\delta Q_{\rm rev}$，且热量足够小，使温度在这一微元过程中可视
为不变，则熵的改变为
\begin{equation}
  \boxed{\dd S=\frac{\delta Q_{\rm rev}}{T}}.
\end{equation}
可逆吸热时熵增加，可逆放热时熵减小；相同热量在较低温度下传递，引起的熵变更大。可逆绝热过程没有
热量交换，因而也是等熵过程。与内能相似，经典热力学首先给出的是熵在过程中的改变，而不是任意系统
熵的绝对数值。

不可逆过程同样会导致熵变。由于熵是状态变量，虽然不能把不可逆路径上的每一中间状态都当作平衡态，
仍可另选一条连接同一初、末态的可逆路径来计算 $\Delta S$。对孤立系统，任意自发过程只可能使熵增加；
在理想可逆极限中熵保持不变，孤立系统中的过程不能使总熵减小。这就是熵增原理。其直接推论是：孤立
系统的热平衡态是熵最大的状态。

在熵概念提出的早期，Clausius 曾把这一结论直接应用于整个宇宙，设想宇宙最终达到处处均匀、温度相同
的热平衡状态，所有细节全部消失。今天我们知道，不能把一个简单孤立热力学系统的结论不加条件地套用
到演化的宇宙；宇宙中丰富结构的形成还受到引力、膨胀和非平衡过程等因素控制。

\subsubsection*{准静态为什么还不够}

准静态只表示变化足够缓慢，使系统能用一串近似平衡态描述；可逆还要求系统和环境都能在不留下
其他变化的条件下复原。例如，缓慢拖动有摩擦的活塞可以近似准静态，但摩擦仍把机械能转为内能，
单纯反向拖动不能把这些内能全部恢复为原来的机械能。有限温差传热也有同样的问题。
因此，可逆传热须令两侧温差趋于零，可逆体积功须令内外压强差趋于零，并排除摩擦等耗散。
这是一种极限模型，不是要求实际装置在有限时间内严格达到它。

\subsubsection*{熵流与熵产生}

对不交换物质的系统，热量可以从不同温度的环境输入。设第 $j$ 个恒温热库温度为 $T_j$，
在全过程向系统传入的带符号热量为 $Q_j$。热库自身的熵变为 $-Q_j/T_j$。
把系统和这些热库一起考察，第二定律要求总熵变非负，因而
\begin{equation}
 \Delta S_{\rm sys}-\sum_j\frac{Q_j}{T_j}=S_{\rm gen}\ge0.
\end{equation}
这里 $S_{\rm gen}$ 是系统连同有关传热边界的总熵产生；它不是新的能量。
等号对应整个过程可逆。如果环境温度变化，应把 $Q_j/T_j$ 换成相应的积分，不能随意选一个平均温度。
对循环，$\Delta S_{\rm sys}=0$，于是
\begin{equation}
 \sum_j\frac{Q_j}{T_j}\le0.
\end{equation}
这就是以恒温热库分段表示的克劳修斯不等式。循环可以把吸热和放热的熵交换相互抵消，
却不能用负的熵产生去抵消耗散。

例如，两热库间直接传递热量 $Q>0$，高温为 $T_h$、低温为 $T_c<T_h$。
高温热库熵变 $-Q/T_h$，低温热库熵变 $Q/T_c$，故
\[
 S_{\rm gen}=Q\left(\frac1{T_c}-\frac1{T_h}\right)>0.
\]
总能量没有损失，但可以利用的温差减小了。若 $T_h\to T_c$，单位热量引起的熵产生趋于零，
这正是可逆传热必须取无穷小温差的原因。

\subsubsection*{由熵不等式判断热机与制冷机}

在 $T_h>T_c>0$ 的两个热库之间，热机每循环从高温热库取热 $Q_h>0$，向低温热库排热 $Q_c>0$。
工质循环复原，环境总熵变为
\[
 -\frac{Q_h}{T_h}+\frac{Q_c}{T_c}\ge0
 \quad\Longrightarrow\quad
 \frac{Q_c}{Q_h}\ge\frac{T_c}{T_h}.
\]
结合 $W=Q_h-Q_c$，得到
\begin{equation}
 \eta=1-\frac{Q_c}{Q_h}\le1-\frac{T_c}{T_h}.
\end{equation}
这次推导没有使用气体物态方程，所以卡诺效率是任意工质的上限，而不只是理想气体计算的巧合。
若装置的某一部分不可逆，通常有严格小于号。

制冷机的有用效果是从低温热库取出 $Q_c$，外界输入功 $W>0$，向高温热库排出
$Q_h=Q_c+W$。此时热库熵变满足
\[
 \frac{Q_h}{T_h}-\frac{Q_c}{T_c}\ge0
 \quad\Longrightarrow\quad
 W\ge Q_c\frac{T_h-T_c}{T_c}.
\]
制冷系数和供热系数分别为
\begin{equation}
 \varepsilon_{\rm ref}=\frac{Q_c}{W}\le\frac{T_c}{T_h-T_c},\qquad
 \varepsilon_{\rm hp}=\frac{Q_h}{W}=1+\varepsilon_{\rm ref}
 \le\frac{T_h}{T_h-T_c}.
\end{equation}
二者可以大于1，因为它们比较的是搬运热量与输入功，而不是把输入热量转为功的比例。
例如 $T_h=300\,\mathrm K,T_c=270\,\mathrm K$ 时，可逆制冷系数为9；搬走
$900\,\mathrm J$ 的热量至少需要 $100\,\mathrm J$ 的功，向高温环境排热为
$1000\,\mathrm J$。由 $1000/300=900/270$ 可同时核验熵平衡。


\subsection*{具有简单物态方程系统的熵}

在绝大多数情形下，热力学系统的熵不容易直接求出；但对物态较简单的系统，某些过程中的熵变可以精确
或近似计算。最简单的系统之一是热库。它有极大的热容量，无论吸收还是放出所研究的有限热量，其温度
都可以近似认为不变。因此，温度为 $T$ 的热库吸收热量 $Q$ 时熵变就是 $Q/T$。

另一类系统以固体和液体为代表。通常吸、放热会引起温度变化，虽然物质体积也随温度改变，但固体、液体
的体积变化往往很小。若忽略体积变化，就同时忽略了系统与环境通过体积功发生的能量交换。设系统比热
为 $c$，并在所研究温区内近似为常数。质量为 $m$ 的物质在温度为 $T$ 时可逆吸收
$\dd Q=mc\,\dd T$，因而
\[
  \dd S=\frac{\dd Q}{T}=mc\frac{\dd T}{T}.
\]
从参考温度 $T_0$ 积分到任意温度 $T$，得到
\begin{equation}
  \boxed{S(T)-S(T_0)=mc\ln\frac{T}{T_0}}.
\end{equation}

理想气体的熵不仅依赖温度，还依赖体积。设有 $n$ 摩尔理想气体，参考状态为 $(V_0,T_0)$，要计算任意
状态 $(V,T)$ 相对参考状态的熵变，先假设 $c_v$ 在温区内为常数。由于熵是状态量，可以任意设计一条连接两状态的可逆路径。图21.5
选择两步路径：先作等容过程 $(V_0,T_0)\to(V_0,T)$，再作等温过程 $(V_0,T)\to(V,T)$。第一步熵变为
\[
  S(V_0,T)-S(V_0,T_0)=n c_v\ln\frac{T}{T_0};
\]
第二步中理想气体的吸热量等于对外功，故
\[
  S(V,T)-S(V_0,T)=nR\ln\frac{V}{V_0}.
\]
两式相加即得
\begin{equation}
  \boxed{S(V,T)-S(V_0,T_0)
  =n c_v\ln\frac{T}{T_0}+nR\ln\frac{V}{V_0}}.
\end{equation}

\entropyPathFigure
\diagramcaption{图 21.5：理想气体任意两状态间的熵差。图中用等容升温加等温膨胀的可逆路径代替任意真实路径。}
图21.5的 $A\to B$ 改变温度而不改体积，$B\to C$ 改变体积而不改温度；因此两个对数项各有清楚的来源。
若实际过程是降温或压缩，相应积分上限小于下限，对数自动给出负号，并不要求实际过程也按图运行。
若 $c_v$ 随温度变化，应保留积分
\[
 \Delta S=n\int_{T_0}^{T}\frac{c_v(T')}{T'}\,\dd T'+nR\ln\frac{V}{V_0}.
\]

\subsection*{熵的微观解释}

宏观量内能可以解释为微观粒子动能与势能的总和。熵也最初来自宏观热过程的研究，但随着分子运动理论
的发展，人们逐渐认识到它的微观含义。Boltzmann 最早系统提出基于分子运动理论的熵解释。虽然这一理论
起初饱受争议，但后来被广泛接受；在主流学界尚未普遍接受 Boltzmann 理论时，Planck 和 Einstein 已经
基于其统计思想建立早期量子理论。

\boltzmannMemorialFigure
\diagramcaption{图 21.6：玻尔兹曼纪念碑上的熵公式，把宏观熵与微观状态的可能数目联系起来。}
图21.6中纪念碑上方的公式，是接下来计数解释的核心。

今天通常把 Boltzmann 公式写成
\begin{equation}
  \boxed{S=k_B\ln\Omega},
\end{equation}
其中 $k_B$ 是 Boltzmann 常数，$\Omega$ 是给定宏观状态所对应的微观状态数。一个宏观状态允许的微观
排布可能性越多，熵就越大；这与把熵理解为系统“无序程度”的传统图像接近，但微观状态数的说法更准确。

先看气体分布这一例子。把容器分为等大的左、右两半，暂只记录每个粒子位于哪一半，而不记录具体位置和速度。
给 $N$ 个粒子临时编号，是为了计数各粒子的左右选择，不是把编号当作新的物理自由度。
若左侧有 $N_1$ 个，右侧有 $N-N_1$ 个，则先从全部粒子中挑出左侧粒子的有序序列，有
$N(N-1)\cdots(N-N_1+1)$ 种；同一组左侧粒子的内部排列被重复计了 $N_1!$ 次，故应除去：
\[
 \Omega_{\rm LR}(N_1)=\frac{N(N-1)\cdots(N-N_1+1)}{N_1!}
 =\frac{N!}{N_1!(N-N_1)!}.
\]
例如 $N=4$ 时，全在左侧仅1种，两边各2个有6种。一般有
$\Omega_{\rm LR}(N_1+1)/\Omega_{\rm LR}(N_1)=(N-N_1)/(N_1+1)$，
所以计数先增后减，在左右数目最接近时最大。每个粒子左右等可能且彼此独立时，
概率为 $\Omega_{\rm LR}/2^N$。这解释了为什么宏观气体倾向于分散到整个容器。
$\Omega_{\rm LR}$ 只计算这一粗粒化的分布因素，并不是包含全部位置、动量的完整微观状态数。

同一组合计数也可用于顺磁性物质。某些原子具有非零磁矩，其性质有点像指南针；在外磁场作用下，简化模型中每个磁矩
只有沿磁场或逆磁场两个方向。设系统含 $N$ 个原子，其中 $N_1$ 个磁矩向上，剩下
$N_2=N-N_1$ 个向下。宏观测量只能确认上下数目的差，不能确定究竟哪一些原子向上。于是，在 $N$ 个
原子中选出 $N_1$ 个向上的微观排布数为
\begin{equation}
  \boxed{\Omega=\frac{N!}{N_1!(N-N_1)!}}.
\end{equation}
在没有外场时，$N_1=N/2$ 附近的微观状态数最大，这正是通常情况下顺磁性物质不显示宏观磁性的统计
原因。

中学阶段通常不会进行复杂的微观熵运算，但必须注意微观世界与宏观世界之间巨大的数目差异。设一个人
吃下一份标称化学能为 $\SI{40}{J}$ 的食物，体温为 $\SI{310}{K}$，能否把人体熵增算成 $40/310$？
不能：进食是物质流入和化学反应，不等于向恒温封闭系统传入40焦耳的热量；还要考虑物质携熵、散热和代谢。
为比较数目大小，改考察一个明确的理想模型：温度恒为 $\SI{310}{K}$ 的热库净吸收
$\SI{40}{J}$ 热量。这个热库的熵变才是
\[
  \Delta S=\frac{Q}{T}
  =\frac{\SI{40}{J}}{\SI{310}{K}}
  \approx\SI{0.13}{J.K^{-1}}.
\]
另一方面，Boltzmann 公式给出
\[
  \Delta S=k_B\ln\frac{\Omega_f}{\Omega_i}.
\]
联立两式，取 $k_B=\SI{1.380649e-23}{J.K^{-1}}$，得到
\begin{equation}
  \boxed{\ln\frac{\Omega_f}{\Omega_i}
  =\frac{40/310}{1.380649\times10^{-23}}
  \approx9.35\times10^{21}}.
\end{equation}
这里先报告对数，避免把数量级粗略相近的巨大指数误写成状态数之比的数值近似。按这一统计描述，
该比值可写为 $\exp(9.35\times10^{21})$ 的量级说明；宏观上很小的热量也可对应极大的微观可能数目变化。
